\documentclass[leqno]{article}
\usepackage{verbatim}
\usepackage{array}
\usepackage{listings}
\usepackage{fancyvrb}
\usepackage{enumitem}

\usepackage{lmodern}

\usepackage[utf8]{inputenc}
\usepackage[T1]{fontenc}
\usepackage{textcomp}
\usepackage{multicol} \usepackage{mathtools}
\usepackage{amsmath}
\usepackage{wrapfig}
\usepackage{amssymb}
\usepackage{extpfeil}
\usepackage{amsmath,amsfonts,amssymb,amsthm,epsfig,epstopdf,titling,url,array}
\usepackage{hyperref}
\usepackage{eso-pic}
\usepackage{pgf}
\usepackage{tikz}
\usepackage{tikz-cd}
\usepackage{graphicx}

% figure support
\usepackage{import}
\usepackage{xifthen}
\pdfminorversion=7
\usepackage{pdfpages}
\usepackage{transparent}
\usepackage{xcolor}

% geometry
\usepackage{geometry}
\geometry{a4paper, margin=1in}

% paragraph length
\setlength{\parindent}{0em}
\setlength{\parskip}{1em}

\newtheorem*{theorem}{Theorem}
\newtheorem*{lemma}{Lemma}
\newtheorem*{proposition}{Proposition}
\newtheorem*{definition}{Definition}
\newtheorem*{observation}{Observation}
\newtheorem*{example}{Example}

\DeclareMathOperator{\Hom}{Hom}
\DeclareMathOperator{\Poset}{Poset}
\DeclareMathOperator{\im}{Im}
\DeclareMathOperator{\coker}{Coker}
\DeclareMathOperator{\coim}{Coim}

\newcommand{\com}[1]{\textcolor{red}{#1}}
\newcommand{\incfig}[1]{%
\center
\def\svgwidth{0.9\columnwidth}
\import{./figures/}{#1.pdf_tex}
}
\newcommand{\incimg}[1]{%
\center
\includegraphics[width=0.9\columnwidth]{images/#1}
}
\pdfsuppresswarningpagegroup=1


\begin{document}

% Lecture 2

% Lecture 3

\begin{definition}[Diagramatic invariant] $I : \text{knot} \to S$ with $S$ set such that $I(K_1) = I(K_2)$ if $K_1$ and $K_2$ are isotopic.
\end{definition}

\begin{example} The crossing at a cross $ \varepsilon (c) \in \{+1, -1\} $. The linking number of a link $L = K_1 \cup K_2$ is an invariant defined as 
\[
\text{lk}(K_1, K_2) = \frac{1}{2} \sum_{c \in K_1 \cap K_2} \varepsilon(c).
\]
\end{example}

\begin{definition}[Fox $p$-coloring] Let $D$ a knot diagram with arcs $x_1, \ldots, x_m$. A Fox $p$-coloring of $D$ is an assignment of colors
\[
x_i \mapsto a_i \in \mathbb{Z}_p
\]
such that at each crossing, we have $a, c$ underarcs and $b$ overarc, then $2b \equiv a + c \pmod{p}$.

We have a homogeneous system of equations $C_D x = 0$ 

\end{definition}

\begin{proposition} The dimension of 
    \[
    \text{Col}_p(D) = \ker  (C_D)
    \]
(hence, the number of Fox $p$-colorings) is a knot invariant.
\end{proposition}

\begin{definition}[Knot group] Let $K$ be a knot. The knot group of $K$ is $\pi_1(S^3 \setminus K)$.
\end{definition}

\begin{definition} $N(K)$ tubular nbh of $K$.
\end{definition}

\begin{definition}[Meridian and longitude] They generate $\pi _1(\partial N(K)) \cong \mathbb{Z}^2$.
\end{definition}

\begin{theorem}[Wirtinger presentation] $g_i$ generator for each arc $x_i$. Then the knot group has presentation
\[
G_K = \langle g_1, \ldots, g_m \mid r_1, \ldots, r_n \rangle
\]
where $r_j : g_k = g_i g_j g_i^{-1}$ are wirtinger relations at each crossing.
\end{theorem}


% TODO: dihedral group

\begin{definition} $C_D^{i,j}$ deleting the $i$-th row and $j$-th column of $C_D$. Then
\[
\det (K) = \left| \det (C_D^{i,j}) \right|  = \left| \Delta_K(-1) \right| = \left| H_1(\Sigma_2(K)) \right|
\]
\end{definition}

\begin{example} Trifoil
\[
C_D = \begin{bmatrix}2 & -1 & -1 \\
-1 & 2 & -1 \\
-1 & -1 & 2 \end{bmatrix}
\]
\end{example}

\begin{theorem} $K$ admits a non-trivial Fox $p$-coloring if and only if $p \mid \det (K)$.
\end{theorem}

\begin{definition}[Writhe] Not a knot invariant (fails Redemeister I)
\[
w(D) = \sum_{c} \varepsilon(c)
\]
\end{definition}

\begin{definition}[Framed knot] A knot $K$ equiped with $v(x)$ a normal vector field. Self-linking number
\[
\text{sl}(K) = \text{lk}(K, K_v = \{x + \varepsilon v(x) : x \in K\})
\]
\end{definition}

% Lecture 4
\section{Seifert and Alexander}

\begin{definition}[Seifert surface] For a knot $K$, a Seifert surface is a compact connected orientable surface $F$ such that $\partial F = K$.
\end{definition}

\begin{proposition} Seifert's algorithm:
TODO
\end{proposition}

\begin{definition}[Euler characteristic]
\[
\chi(F) = s(D) - c(D) = 1 - 2g(F)
\]
where $s(D)$ is the number of Seifert circles and $c(D)$ is the number of crossings.
\end{definition}

\begin{proposition} $K_1\# K_2$ induces $F_1\natural F_2$ and 
\[
g(F_1 \natural F_2) = g(F_1) + g(F_2), \quad g(K_1 \# K_2) = g(K_1) + g(K_2).
\]
\end{proposition}

\begin{proposition} $F$ Seifert surface of $K$ implies $H_1(F, \mathbb{Z}) \cong \mathbb{Z}^{2g(F)}$
\end{proposition}

\begin{definition}[Seifert matrix] Basis $\alpha _1, \ldots, \alpha_{2g}$ of $H_1(F, \mathbb{Z})$, the Seifert matrix $V$ is defined by
\[
V_{ij} = \text{lk}(\alpha_i, \alpha_j^+)
\]
where $\alpha_j^+$ is a push-off of $\alpha_j$ in the positive normal direction of $F$.
\end{definition}

\begin{definition}[Algebraic intersection number] $F$ oriented, $a, b$ oriented simple closed curves on $F$. Then
\[
a \cdot b = \sum_{p \in a \cap b} \varepsilon_p(a, b), \quad \varepsilon_p(a, b) = \begin{cases} +1 & \text{if } (a, b) \text{ is + oriented at } p \\ -1 & \text{if } (a, b) \text{ is - oriented at } p \end{cases}
\]
we also see that $\alpha_i \cdot \alpha_j = V_{ij} - V_{ji}$
\end{definition}

\begin{definition}[Alexander polynomial] $V$ Seifert matrix for $K$. 
    \[
\Delta _K(t) \coloneqq  \det (V - t V^T) \in \mathbb{Z}[t, t^{-1}], \quad \Delta_K(t) = \Delta_K(t^{-1}), \quad \Delta_K(1) = 1.
    \]
\end{definition}

\begin{example} Trefoil $V = \begin{bmatrix} 1 & 0 \\ -1 & 1 \end{bmatrix}$, $\Delta_K(t) = t^2 - t + 1 = t-1+t^{-1}$
\end{example}





\end{document}